\subsection*{Question 7: Platform mass and effective heave mass}

The platform mass is found from hydrostatic equilibrium, i.e. the mass of the
platform is equal to the mass of displaced seawater. The density of seawater is
taken as
\[
    \rho = 1025~\mathrm{kg/m^3}.
\]

From Appendix A, the main dimensions are
\[
\begin{aligned}
    T &= 18.5~\mathrm{m}, &
    L &= 72.3~\mathrm{m}, &
    D &= 11.3~\mathrm{m},\\
    a &= 30.5~\mathrm{m}, &
    b &= 5.40~\mathrm{m}, &
    f &= T - 2b = 7.70~\mathrm{m}.
\end{aligned}
\]
The column radius is
\[
    r_c = \frac{D}{2} = 5.65~\mathrm{m},
\]
and the cross-sectional area of one column is therefore
\[
    A_c = \pi r_c^2 = 100.287~\mathrm{m^2}.
\]

The displaced volume from the four columns is calculated using the submerged
column length down to the top of the pontoons,
\[
    V_{\mathrm{columns}} = 4 A_c f
    = 4 \pi \left(\frac{D}{2}\right)^2 f
    = 3088.855~\mathrm{m^3}.
\]
The pontoons are approximated as rectangular boxes with cross-section
\((2b)\times(2b)\). From the double-symmetric top view, two pontoons are
full length \(L\), while the two side pontoons run between the columns and
therefore have length \(L-2D\). Thus,
\[
    V_{\mathrm{pontoons}}
    = \left[2L + 2(L-2D)\right](2b)(2b)
    = 28460.160~\mathrm{m^3}.
\]
Hence, the total displaced volume is
\[
    \nabla
    = V_{\mathrm{columns}} + V_{\mathrm{pontoons}}
    = 31549.015~\mathrm{m^3}.
\]
The platform mass becomes
\[
    M = \rho \nabla
    = 3.2338 \times 10^7~\mathrm{kg}
    = 32337.740~\mathrm{tons}.
\]

The hydrodynamic added mass is estimated from Appendix B. For the circular
columns in vertical motion,
\[
    C_{m,c} = 1.00,
    \qquad
    A_{R,c} = \pi r_c^2.
\]
Thus,
\[
    A_{33,c}
    = 4 C_{m,c} \rho A_{R,c} f
    = 3.1661 \times 10^6~\mathrm{kg}
    = 3166.076~\mathrm{tons}.
\]

For the pontoons, the section is approximated as square, with \(a/b = 1.00\).
From Appendix B,
\[
    C_{m,p} = 1.51,
    \qquad
    A_{R,p} = \pi b^2.
\]
The pontoon added mass is therefore
\[
    A_{33,p}
    = C_{m,p} \rho A_{R,p}\left[2L + 2(L-2D)\right]
    = 3.4596 \times 10^7~\mathrm{kg}
    = 34596.171~\mathrm{tons}.
\]

The total added mass in heave is
\[
    A_{33} = A_{33,c} + A_{33,p}
    = 3.7762 \times 10^7~\mathrm{kg}
    = 37762.247~\mathrm{tons}.
\]
Finally, the effective mass in heave is
\[
    M_{\mathrm{eff}} = M + A_{33}
    = 7.0100 \times 10^7~\mathrm{kg}
    = 70099.987~\mathrm{tons}.
\]

\subsection*{Question 8: Equation of uncoupled and undamped heave motion}

The heave motion is treated as a single-degree-of-freedom motion. Coupling to all
other rigid-body modes is neglected, and damping is also neglected. With \(z(t)\)
as the heave displacement, positive upwards, the equation of motion is
\[
    \left(M + A_{33}\right)\ddot{z}(t) + C_{33} z(t) = F_3(t),
\]
where \(F_3(t)\) is the vertical wave excitation force.

The hydrostatic restoring coefficient in heave is
\[
    C_{33} = \rho g A_{\mathrm{wp}},
\]
where the waterplane area is due to the four surface-piercing columns,
\[
    A_{\mathrm{wp}}
    = 4 \pi \left(\frac{D}{2}\right)^2
    = 401.150~\mathrm{m^2}.
\]
Therefore,
\[
    C_{33}
    = 1025 \cdot 9.81 \cdot 401.150
    = 4.0337 \times 10^6~\mathrm{N/m}.
\]

For harmonic excitation at angular frequency \(\omega\), the frequency-domain
form of the same equation is
\[
    \left[C_{33} - \omega^2\left(M + A_{33}\right)\right]\hat{z}
    = \hat{F}_3(\omega).
\]
This form is used later to determine the heave response amplitude operator.
